\chapter{Forecasting Period} \label{chap:forecasting_period}

The Discounted Cash Flow model implicitly assumes that beyond a certain forecast horizon, a company reaches a "steady-state." This conceptual steady-state firm is characterized by predictable performance metrics, including stable growth rates, profit margins, reinvestment needs (CAPEX and working capital), and cost of capital (reflected in the WACC - Weighted Average Cost of Capital).

While companies in steady state don't necessarily share identical parameters, they exhibit significantly higher predictability in their operations and financial trajectory compared to firms in transitional phases.

The forecasting period, \(N\), represents the number of years for which future cash flows are estimated (as discussed in Chapter \ref{chap:cash_flow}) before entering the steady-state assumption. The transition to the steady-state at year, \(N\), occurs for primarily two reasons:

\begin{enumerate}
    \item \textbf{Convergence to Steady State}: The company's key financial characteristics (e.g., growth, margins, reinvestment rates, cost of capital) have converged to stable, long-term sustainable levels
    \item  \textbf{Forecast Reliability Limit}: Projecting detailed, distinct annual cash flows beyond year \(N\) becomes increasingly speculative and impractical due to inherent uncertainties. Relying on a simplified steady-state model becomes more robust than making highly uncertain long-range forecasts.
\end{enumerate}

For some highly stable, mature companies (as Coca-Cola or Procter \& Gamble), their current operations may already closely resemble a steady state. Their year-over-year performance changes little, suggesting convergence has largely already occurred. In such cases, making detailed multi-year forecasts offers limited additional insight compared to a steady-state assumption, making explicit forecasting potentially unnecessary or impractical. Consequently, \(N\) might be set to \(0\), and the valuation proceeds directly using a terminal value calculation based on current or next-year expected performance.

For most companies, however, a non-zero forecast period \(N\) is required to capture anticipated changes before stability is reached. The appropriate length of \(N\) depends on the analyst's judgment regarding how long significant transitional dynamics (high growth, margin changes, major investments) are expected to persist. Companies with greater visibility into their future and operating in more stable industries might warrant a longer forecast period \(N\) where explicit estimates remain meaningful. Conversely, companies in highly volatile industries or facing significant uncertainty might necessitate a shorter \(N\) before resorting to the steady-state assumption.

While the choice of \(N\) involves analyst judgment, common practice often limits \(N\) to a maximum of \(10\) years. This convention helps mitigate overconfidence in long-range forecasting and maintains model integrity. Although an upper limit is often advised, there isn't a prescribed lower limit when \(N>0\). The key is determining the point beyond which a detailed annual forecast ceases to add value compared to the steady-state simplification \parencite{damodaran2025investment, koller2020valuation}.
